\section{Original FFT Results}
\label{sec:legacy-results}
\subsection{Numerical agreement and completion}
All 1320 incremental-versus-complete schedule cases passed, with bitwise-identical
completed spectra within each tested type. The independent reference checked
6048 complex bins across both precisions. The maximum scaled errors were
$8.69\times10^{-8}$ for float (rounded upward) and
$1.06\times10^{-14}$ for double (rounded upward). These measurements support
correctness for the tested fixtures, sizes, and windows; they do not replace an
error analysis of every supported type or the module's amplitude calibration.

The schedule tests confirm Table~\ref{tab:cadence} and the distinction between
``no later than $H$'' and ``exactly at $H$.'' Increasing the requested horizon
can leave the actual completion interval unchanged until the integer quota
changes. Thus frame rate is a staircase function of the horizon under
completion-driven restarts, not generally its reciprocal.

\subsection{Peak work placement}
Table~\ref{tab:timing} reports the $H=N/2$ subset; the artifact also contains
$H=N/4$. Figure~\ref{fig:peaks} includes the full range of observed maxima across
repetitions. The median of run maxima is much smaller in incremental mode,
but individual maxima vary considerably. At $N=2048$, for example, one
incremental run reaches 21.209 microseconds even though its median run maximum
is 1.625 microseconds. Such observations are consistent with an ordinary
scheduled process and are retained rather than removed as outliers.

\begin{table}[tb]
\centering\small
\begin{tabular}{@{}rrr rrr@{}}
\toprule
& \multicolumn{2}{c}{Frame time ($\mu$s)} & \multicolumn{2}{c}{Run maximum ($\mu$s)} & Peak ratio \\
\cmidrule(lr){2-3}\cmidrule(lr){4-5}
$N$ & Complete & Incremental & Complete & Incremental & Paired 95\% interval \\
\midrule
1,024 & 7.619 & 7.518 & 7.917 & 0.875 & [0.083, 0.508] \\
2,048 & 16.662 & 15.915 & 17.291 & 1.625 & [0.091, 0.111] \\
4,096 & 36.620 & 35.603 & 39.458 & 3.333 & [0.079, 0.296] \\
16,384 & 169.510 & 163.844 & 179.625 & 13.125 & [0.073, 0.113] \\
\bottomrule
\end{tabular}
\caption{Medians across nine repetitions, $H=N/2$. Frame time comes from the aggregate pass; run maximum comes from the separately instrumented pass. Intervals are bootstrap intervals for the median paired incremental/complete run-maximum ratio, not ratios of the displayed marginal medians and not worst-case bounds.}
\label{tab:timing}
\end{table}


\begin{figure}[tb]
\centering
\begin{tikzpicture}
\begin{axis}[width=0.93\linewidth,height=6.3cm,ymode=log,xmode=log,log basis x=2,
 ylabel={Observed run maximum ($\mu$s)},xlabel={Real transform length $N$},
 xtick={1024,2048,4096,16384},xticklabels={1024,2048,4096,16384},xmin=800,xmax=21000,
 ymajorgrids=true,grid style={black!10},legend pos=north west,
 legend style={draw=none,font=\small},tick label style={font=\small}]
\addplot+[color=ink,mark=square*,error bars/.cd,y dir=both,y explicit]
 coordinates {(1024,7.917) += (0,14.374) -= (0,0.167)
 (2048,17.291) += (0,6.084) -= (0,0.457)
 (4096,39.458) += (0,6.751) -= (0,1.958)
 (16384,179.625) += (0,13.084) -= (0,7.959)};
\addlegendentry{Complete}
\addplot+[color=accent,mark=*,error bars/.cd,y dir=both,y explicit]
 coordinates {(1024,0.875) += (0,5.584) -= (0,0.083)
 (2048,1.625) += (0,19.584) -= (0,0.083)
 (4096,3.333) += (0,14.333) -= (0,0.249)
 (16384,13.125) += (0,12.333) -= (0,0.542)};
\addlegendentry{Incremental}
\end{axis}
\end{tikzpicture}
\caption{Median and full min--max range of per-run maximum call duration over nine repetitions, $H=N/2$. Vertical bars are observed ranges, not confidence intervals. Both axes are logarithmic; horizontal ticks identify the actual transform lengths. Measurements include RFFT preparation and reconstruction, but exclude module smoothing and graphics preparation.}
\label{fig:peaks}
\end{figure}


The incremental preparation-call medians are 0.542, 1.125, 2.709, and
10.667 microseconds for the four lengths. Completing-call medians are 0.750,
1.500, 3.042, and 12.375 microseconds. Their growth is consistent with the bulk
work identified by source analysis, although timings alone do not prove an
asymptotic bound. Distributing only the butterfly traversal leaves both ends of
the transform visible in the cost profile.

A striking measurement pitfall is that the median per-run p99 is approximately
42 ns for \emph{both} modes at all four lengths in this subset. The complete
mode's expensive call occurs only once per $H$ calls, below the top one percent
for these horizons. A p99-only comparison would largely miss the burst being
studied. Phase-specific measurements and run maxima are needed alongside
quantiles; none of them by itself is a proven worst-case execution time.

\subsection{Aggregate work and interpretation}
The median paired incremental/complete frame-time ratios for increasing $N$
are approximately 0.979, 0.955, 0.972, and 0.967. The corresponding percentile
intervals are [0.973, 1.005], [0.955, 0.956], [0.960, 0.977], and
[0.963, 0.968]. These modest differences are specific to the two loop structures
as optimized by this compiler on this workload. They do not indicate a reduction
in the butterfly count, nor establish an algorithmic throughput advantage. A
different compiler, batching policy, cache state, or vectorized library could
reverse the result. The robust interpretation of this campaign is the change
in work placement within this implementation, with broadly comparable total
frame cost.

This original campaign uses repeated fixed inputs, warm tables, and a scalar standalone
process. They omit ring insertion, Rack parameter access, SIMD channels,
frequency/time smoothing, display conversion, GUI consumption, and other
modules. They therefore cannot establish whole-patch dropout rates or a maximum
number of safe analyzer instances. Neither timing campaign measures a running Rack session.

